
\documentclass{article}
\usepackage{colortbl}
\usepackage{makecell}
\usepackage{multirow}
\usepackage{supertabular}

\begin{document}

\newcounter{utterance}

\centering \large Interaction Transcript for game `cladder', experiment `full\_v1.5\_default', episode 8689 with reflexion{-}Qwen3.5{-}9B.
\vspace{24pt}

{ \footnotesize  \setcounter{utterance}{1}
\setlength{\tabcolsep}{0pt}
\begin{supertabular}{c@{$\;$}|p{.15\linewidth}@{}p{.15\linewidth}p{.15\linewidth}p{.15\linewidth}p{.15\linewidth}p{.15\linewidth}}
    \# & \multicolumn{2}{c}{Player} && \multicolumn{2}{c}{Game Master} \\
    \hline

    \theutterance \stepcounter{utterance}  
    & & & \multicolumn{4}{p{0.6\linewidth}}{
        \cellcolor[rgb]{0.9,0.9,0.9}{
            \makecell[{{p{\linewidth}}}]{
                \texttt{\tiny{[P1$\langle$GM]}}
                \texttt{You are an expert in causal inference. The following question is not a typical commonsense query, but rather a meticulously designed question created by a professor specializing in causal inference, intended to assess the students' mastery of the course content.} \\
\texttt{Imagine a self{-}contained, hypothetical world with only the following conditions, and without any unmentioned factors or causal relationships: Jyka has a direct effect on qwiu and yupt. Yupt has a direct effect on kwox. Qwiu has a direct effect on kwox. The overall probability of jyka is 53\%. For those who are not jyka, the probability of kwox is 60\%. For those who are jyka, the probability of kwox is 50\%. Is kwox less likely than not kwox overall?} \\
\texttt{Start your answer with 'Yes' or 'No', followed by additional reasoning or evidence to support your explanation.} \\
            }
        }
    }
     \\ \\

    \theutterance \stepcounter{utterance}  
    & \multicolumn{4}{p{0.6\linewidth}}{
        \cellcolor[rgb]{0.561,0.737,0.561}{
            \makecell[{{p{\linewidth}}}]{
                \texttt{\tiny{[P1$\rangle$GM]}}
                \texttt{No} \\
\\ 
\texttt{**Reasoning:**} \\
\\ 
\texttt{To determine if \$kwox\$ is less likely than not \$kwox\$ overall, we must calculate the marginal probability \$P(kwox)\$ and compare it to \$0.5\$.} \\
\\ 
\texttt{1.  **Identify the given probabilities:**} \\
\texttt{    *   \$P(jyka) = 0.53\$} \\
\texttt{    *   \$P(kwox | \textbackslash{}neg jyka) = 0.60\$} \\
\texttt{    *   \$P(kwox | jyka) = 0.50\$} \\
\\ 
\texttt{2.  **Apply the Law of Total Probability:**} \\
\texttt{    The overall probability of \$kwox\$ is the weighted sum of the conditional probabilities based on the status of \$jyka\$:} \\
\texttt{    \$\$P(kwox) = P(kwox | jyka) \textbackslash{}cdot P(jyka) + P(kwox | \textbackslash{}neg jyka) \textbackslash{}cdot P(\textbackslash{}neg jyka)\$\$} \\
\\ 
\texttt{3.  **Calculate the components:**} \\
\texttt{    *   \$P(\textbackslash{}neg jyka) = 1 {-} P(jyka) = 1 {-} 0.53 = 0.47\$} \\
\texttt{    *   Contribution from \$jyka\$: \$0.50 \textbackslash{}times 0.53 = 0.265\$} \\
\texttt{    *   Contribution from \$\textbackslash{}neg jyka\$: \$0.60 \textbackslash{}times 0.47 = 0.282\$} \\
\\ 
\texttt{4.  **Sum the contributions:**} \\
\texttt{    \$\$P(kwox) = 0.265 + 0.282 = 0.547\$\$} \\
\\ 
\texttt{5.  **Compare to the threshold:**} \\
\texttt{    The question asks if \$kwox\$ is *less likely* than *not \$kwox\$*.} \\
\texttt{    *   \$P(kwox) = 0.547\$} \\
\texttt{    *   \$P(\textbackslash{}neg kwox) = 1 {-} 0.547 = 0.453\$} \\
\\ 
\texttt{    Since \$0.547 > 0.453\$, \$kwox\$ is more likely than not \$kwox\$. Therefore, the statement that \$kwox\$ is less likely is false.} \\
\\ 
\texttt{*Note: The information regarding the direct effects of \$jyka\$ on \$qwiu\$ and \$yupt\$, and \$yupt\$ on \$kwox\$, describes the causal graph structure but is not required to answer the specific probabilistic question, as the conditional probabilities for \$kwox\$ given \$jyka\$ are explicitly provided.*} \\
            }
        }
    }
    & & \\ \\

    \theutterance \stepcounter{utterance}  
    & & & \multicolumn{2}{p{0.3\linewidth}}{
        \cellcolor[rgb]{0.9,0.9,0.9}{
            \makecell[{{p{\linewidth}}}]{
                \texttt{\tiny{[GM$|$GM]}}
                \texttt{No} \\
            }
        }
    }
    & & \\ \\

    \theutterance \stepcounter{utterance}  
    & & & \multicolumn{2}{p{0.3\linewidth}}{
        \cellcolor[rgb]{0.9,0.9,0.9}{
            \makecell[{{p{\linewidth}}}]{
                \texttt{\tiny{[GM$|$GM]}}
                \texttt{no} \\
            }
        }
    }
    & & \\ \\

    \theutterance \stepcounter{utterance}  
    & & & \multicolumn{2}{p{0.3\linewidth}}{
        \cellcolor[rgb]{0.9,0.9,0.9}{
            \makecell[{{p{\linewidth}}}]{
                \texttt{\tiny{[GM$|$GM]}}
                \texttt{game\_result = WIN} \\
            }
        }
    }
    & & \\ \\

\end{supertabular}
}

\end{document}
